%! AP Statistics — Unit 1 Cover Sheet
\documentclass[10pt]{article}
\usepackage{apstats-article}
\usepackage{apstats-boxes}

%=============================================================================
\begin{document}

% ── Full-bleed navy banner ────────────────────────────────────────────────────
\begin{tikzpicture}[remember picture, overlay]
  \fill[navy] (current page.north west)
    rectangle ([yshift=-1.4in]current page.north east);
\end{tikzpicture}
\vspace{-1.3in}

\begin{center}
  {\color{white}\bfseries\fontsize{28}{34}\selectfont AP Statistics}\\[8pt]
  {\color{goldacc}\bfseries\Large Shepherd \quad 2026--2027}\\[14pt]
  {\color{white}\bfseries\LARGE Unit 1: Exploring One-Variable Data}\\[4pt]
\end{center}

\vspace{0.30in}

% ── Unit overview ─────────────────────────────────────────────────────────────
\begin{tcolorbox}[colback=sky,colframe=navy,boxrule=0.9pt,arc=2mm,
    left=4mm,right=4mm,top=3mm,bottom=3mm]
  \textbf{Unit Overview.}\quad
  Unit 1 introduces the foundational language and tools of statistics. Students learn
  to classify variables, organize and represent categorical and quantitative data,
  and describe distributions using shape, center, spread, and unusual features.
  The unit culminates with summary statistics, boxplots, comparisons across groups,
  and an introduction to the Normal distribution and standardized scores.
\end{tcolorbox}

\vspace{0.22in}

% ── Lessons in this unit ──────────────────────────────────────────────────────
\renewcommand{\arraystretch}{1.6}
\begin{skillbox}[Lessons in This Unit]{goldbox}
\begin{tabularx}{\linewidth}{@{} c >{\bfseries}p{0.38\linewidth} X @{}}
  \textbf{\#} & \textbf{Title} & \textbf{Focus} \\
  \hline
  1.1 & Introducing Statistics
    & What can we learn from data? Statistical questions, individuals, variables,
      and the data cycle. \\
  \hline
  1.2 & The Language of Variation
    & Variables vs.\ values; categorical vs.\ quantitative variables; distribution. \\
  \hline
  1.3 & Representing a Categorical Variable with Tables
    & Frequency tables, relative frequency tables; marginal and joint frequencies. \\
  \hline
  1.4 & Representing a Categorical Variable with Graphs
    & Bar graphs, segmented bar graphs, pie charts; comparing displays. \\
  \hline
  1.5 & Representing a Quantitative Variable with Graphs
    & Dotplots, stemplots, histograms; choosing appropriate displays. \\
  \hline
  1.6 & Describing the Distribution of a Quantitative Variable
    & Shape, center, spread, and unusual features (SOCS); describing in context. \\
  \hline
  1.7 & Summary Statistics for a Quantitative Variable
    & Mean, median, range, IQR, standard deviation; resistance to outliers. \\
  \hline
  1.8 & Graphical Representations of Summary Statistics
    & Boxplots and modified boxplots; identifying outliers with the 1.5 × IQR rule. \\
  \hline
  1.9 & Comparing Distributions of a Quantitative Variable
    & Parallel boxplots, back-to-back stemplots; comparing shape, center, spread. \\
  \hline
  1.10 & The Normal Distribution and z-Scores
    & Properties of the Normal curve; the 68--95--99.7 rule; standardized scores. \\
  \hline
\end{tabularx}
\end{skillbox}

\vspace{0.18in}

% ── Standards covered ─────────────────────────────────────────────────────────
\begin{skillbox}[AP Statistics Big Ideas \& Learning Objectives — Unit 1]{greenbox}
\begin{tabularx}{\linewidth}{@{} >{\bfseries}p{0.14\linewidth} X @{}}
  DAT-1.A & Identify the individuals and variables in a data set; classify variables as
            categorical or quantitative. \\
  DAT-1.B & Create and interpret frequency tables and graphical displays for categorical data. \\
  DAT-1.C & Create and interpret graphical displays for quantitative data (dotplots,
            stemplots, histograms). \\
  DAT-1.D & Describe and compare distributions using shape, center, spread, and
            unusual features (SOCS). \\
  DAT-1.E & Calculate and interpret measures of center (mean, median) and spread
            (range, IQR, standard deviation) for quantitative data. \\
  DAT-1.F & Construct and interpret boxplots; use the 1.5 × IQR rule to identify outliers. \\
  DAT-1.G & Compare distributions of quantitative data using parallel boxplots or
            back-to-back stemplots. \\
  DAT-1.H & Use the Normal distribution and the 68--95--99.7 rule to estimate
            probabilities; calculate and interpret z-scores. \\
\end{tabularx}
\end{skillbox}

\end{document}
